\section{Introduction}
\label{sec:intro}
% trungnt
Remote terminal access remains a fundamental interface for operating servers, cloud instances, edge devices, and experimental testbeds. While traditionally driven by human operators, recent automation workflows increasingly exercise remote shells through repeated command execution, output parsing, and interactive terminal updates. These automated high-frequency interactions are highly sensitive to transport-level bottlenecks. SSHv2 remains the de facto protocol for remote shell access, but its TCP transport suffers from head-of-line (HoL) blocking and binds sessions to fixed network addresses, causing severe stalls during automated command execution under network impairments \cite{Scharf2007}. Modern alternatives utilizing UDP (Mosh \cite{Winstein2012}) or HTTP/3 and QUIC (SSH3 \cite{michel2023}) address these legacy weaknesses, yet their design focus on human typing and network roaming leaves their reliability for deterministic, burst-driven agent workloads unclear. Therefore, this paper asks: \textbf{Which SSH Connection Protocol can meet the strict latency and completeness demands of closed-loop terminal agents?}

Evaluating the performance of terminal protocols for autonomous agents is inherently challenging because current methodologies isolate either the human experience or the AI logic, but rarely their intersection. In the networking domain, foundational evaluations of protocols like Mosh and SSH3 heavily prioritize human-centric metrics \cite{michel2023}. They measure single-character echo latency and connection resilience during IP roaming, assuming input rates of just a few words per minute. Conversely, recent AI research introduces sophisticated terminal environments to evaluate the reasoning capabilities of AI. However, these AI benchmarks treat the underlying network as a frictionless, zero-latency pipe, focusing exclusively on end-to-end task success \cite{Liu2024}. This dichotomy creates a critical evaluation gap. When a transport protocol stalls due to packet loss or drops terminal output during a high-frequency automated burst, the agent’s closed-loop observation process is fundamentally disrupted. Consequently, existing literature lacks a targeted methodology to measure how different transport mechanisms compromise the completeness and timing of AI-driven interactions. To address this gap, this paper does not evaluate end-to-end AI task success. Instead, we introduce a deterministic benchmark that models an AI-like interactive client. By isolating transport-level behavior, we evaluate the latency and completeness of the feedback channel that closed-loop systems experience over these protocols. 

We design a workload suite that captures three distinct interaction patterns to AI clients: short-output command loops, large-output data streams, and complex single-pane and multi-pane interactive editing sessions. We deploy these workloads across SSHv2, SSH3, and Mosh to quantify per-workload command-response latency and output completeness. To ensure robust validation, we evaluate these metrics across two distinct testing topologies: a controlled testbed featuring three defined network impairment profiles, and a real-world VPN deployment. The main contributions of this paper are summarized as follows:

% The goal is to isolate the transport-level behavior that such closed-loop clients would experience before adding model reasoning, planning, or task-specific policy decisions.
% trungnt


\begin{itemize}
    \item We conduct a deterministic measurement study evaluating SSHv2, SSH3, and Mosh across real-world VPN and controlled network impairments, specifically isolating transport-level bottlenecks for AI-driven terminal workloads.
    
    \item We show that SSH3's QUIC-based architecture is optimal for agile, short-lived agent interactions. It achieves the lowest session setup latency and provides the fastest command-response loop for short terminal outputs.
    
    \item We reveal a tradeoff between human-visible responsiveness and programmatic reliability for large-output commands. While Mosh excels in interactive keystroke latency, its state-synchronization model drops over 98\% of expected bytes, breaking the AI observation loop. Conversely, legacy SSHv2 strictly preserves 100\% byte completeness and completes large data streams significantly faster than SSH3.
\end{itemize}

We discuss background and related work in Section~\ref{sec:background}, detail our methodology in Section~\ref{sec:methodology}, and present our evaluation in Section~\ref{sec:evaluation}, before concluding in Section~\ref{sec:conclusion}.
